\section{Sucesiones y Series Numéricas}

\subsection{Concepto de serie y convergencia}

\begin{definition}{Serie}{def_serie}
    Sea $\{a_n\}_{n\in\mathbb{N}}$ una sucesión. Al sumar los términos de la sucesión se obtiene una serie infinita o simplemente serie y denotaremos:
    \[ \sum_{n=1}^{\infty}a_n = a_1 + a_2 + \dots + a_n + \dots \]
\end{definition}

No confundir una serie con sumatoria: una sumatoria suma una cantidad finita de términos, en cambio una serie suma todos los términos de una sucesión.

Observar que sumar infinitas cantidades no es algo sencillo, ¿podemos garantizar que esta suma existe?, además pueden ocurrir situaciones extrañas, por ejemplo, si queremos conocer la suma de $\displaystyle\sum_{n=1}^{\infty}(-1)^{n+1}$, es decir:
\[ 1 + (-1) + 1 + (-1) + \dots \]
La suma podría ser $1$ ó $0$, dependiendo de cómo se sume. En efecto:
\[ 1 + [(-1) + 1] + [(-1) + 1] + \dots = 1 \]
\[ [1 + (-1)] + [1 + (-1)] + \dots = 0 \]
Para evitar ambigüedades, es necesario precisar qué se entiende como suma de una serie.

\begin{definition}{Sucesión de sumas parciales y convergencia}{def_convergencia}
    A partir de una sucesión $\{a_n\}_{n\in\mathbb{N}}$, podemos formar una nueva sucesión $\{S_N\}_{N\in\mathbb{N}}$, donde:
    \begin{align*}
        S_1 &= a_1 \\
        S_2 &= a_1 + a_2 \\
        S_3 &= a_1 + a_2 + a_3 \\
            &\vdots \\
        S_N &= a_1 + a_2 + \dots + a_N = \sum_{n=1}^{N}a_n
    \end{align*}
    Esta nueva sucesión se conoce como \textbf{sucesión de sumas parciales} de $a_n$.
    Diremos que la serie $\displaystyle\sum_{n=1}^{\infty}a_n$ \textbf{converge} si el límite de la suma parcial $N$-ésima existe:
    \[ \lim_{N\to\infty}S_N = \lim_{N\to\infty}\sum_{n=1}^{N}a_n = \sum_{n=1}^{\infty}a_n = S \in \mathbb{R} \]
    En tal caso, diremos que $S$ es la suma de la serie. En caso contrario la serie diverge.
\end{definition}

\begin{example}{Análisis de convergencia por sumas parciales}{ej:serie_1}
    Analice la convergencia de las siguientes series:
    \begin{multicols}{2}
    \begin{enumerate}
        \item $\displaystyle\sum_{n=1}^{\infty} n$
        \item $\displaystyle\sum_{n=1}^{\infty} \ln\left(\frac{n}{n+1}\right)$
    \end{enumerate}
    \end{multicols}
\end{example}

\begin{soln}
    $\left. \right.$
    \begin{enumerate}
        \item Sumas parciales: $S_1 = 1$, $S_2 = 1+2=3$, $S_3 = 1+2+3=6$, \dots, $S_n = \displaystyle\sum_{i=1}^{n}i = \frac{n(n+1)}{2}$.\\
        Como $\displaystyle\lim_{n\to\infty}S_n = \lim_{n\to\infty}\frac{n(n+1)}{2} = +\infty$, la serie $\displaystyle \sum_{n=1}^{\infty}{n}$ diverge.
        
        \item Sumas parciales:
        \begin{align*}
            S_1 &= \sum_{n=1}^{1}\ln\left(\frac{n}{n+1}\right) = \ln\left(\frac{1}{2}\right) \\
            S_2 &= \sum_{n=1}^{2}\ln\left(\frac{n}{n+1}\right) = \ln\left(\frac{1}{2}\right) + \ln\left(\frac{2}{3}\right) = \ln\left(\frac{1}{2}\cdot\frac{2}{3}\right) = \ln\left(\frac{1}{3}\right) \\
            S_3 &= \sum_{n=1}^{3}\ln\left(\frac{n}{n+1}\right) = \ln\left(\frac{1}{2}\right) + \ln\left(\frac{2}{3}\right) + \ln\left(\frac{3}{4}\right) = \ln\left(\frac{1}{2}\cdot\frac{2}{3}\cdot\frac{3}{4}\right) = \ln\left(\frac{1}{4}\right) \\
            &\vdots \\
            S_N &= \sum_{n=1}^{N}\ln\left(\frac{n}{n+1}\right) = \ln\left(\frac{1}{2}\cdot\frac{2}{3}\cdot\frac{3}{4}\cdots\frac{N}{N+1}\right) = \ln\left(\frac{1}{N+1}\right)
        \end{align*}
        Como $\displaystyle\lim_{N\to\infty}S_N = \lim_{N\to\infty}\ln\left(\frac{1}{N+1}\right) = -\infty$, la serie $\displaystyle\sum_{n=1}^{\infty}\ln\left(\frac{n}{n+1}\right)$ diverge.
    \end{enumerate}
\end{soln}

\begin{proposition}{Propiedades de las series}{prop_series}
    Si $\displaystyle\sum_{n=1}^{\infty}a_n$ y $\displaystyle\sum_{n=1}^{\infty}b_n$ son convergentes y $c \in \mathbb{R}$, entonces:
    \begin{enumerate}
        \item $\displaystyle\sum_{n=1}^{\infty}(a_n \pm b_n) = \sum_{n=1}^{\infty}a_n \pm \sum_{n=1}^{\infty}b_n$
        \item $\displaystyle\sum_{n=1}^{\infty}(c \cdot a_n) = c\sum_{n=1}^{\infty}a_n$
    \end{enumerate}
\end{proposition}

\begin{theorem}{Condición de convergencia}{teo_condicion_convergencia}
    Si $\displaystyle\sum_{n=1}^{\infty}a_n$ converge $\Longrightarrow \displaystyle \lim_{n\to\infty}a_n = 0$.
\end{theorem}

\begin{remark}{El recíproco y el contrarrecíproco}{}
    El teorema solamente nos dice que si la serie converge, entonces el término general de la serie debe tender a cero. Sin embargo, el recíproco no es cierto, es decir, que si el término general de la serie tiende a cero, entonces la serie converge. De hecho, existen series cuyos términos tienden a cero pero divergen, por ejemplo: la serie $\displaystyle \sum_{n=1}^{\infty}{\dfrac{1}{n}}$.
    
    El contrarrecíproco de dicho teorema resulta de mayor utilidad: si el término general de la serie no tiende a cero, entonces la serie diverge.
\end{remark}

\begin{theorem}{Criterio de divergencia}{teo_criterio_divergencia}
    Si $\displaystyle\lim_{n\to\infty}a_n \neq 0 \Longrightarrow \sum_{n=1}^{\infty}a_n$ diverge.
\end{theorem}

Por ejemplo, la serie $\displaystyle\sum_{n=1}^{\infty}e^{\frac{1}{n}}$ diverge, pues $\displaystyle\lim_{n\to\infty}e^{\frac{1}{n}} = 1 \neq 0$.

\subsection{Series notables y descomposiciones}

\subsubsection*{Serie telescópica}
Considere una serie de la forma $\displaystyle\sum_{n=1}^{\infty}(b_n - b_{n+1})$, donde las sumas parciales son:
\begin{align*}
    S_1 &= b_1 - b_2 \\
    S_2 &= (b_1 - b_2) + (b_2 - b_3) = b_1 - b_3 \\
    S_3 &= (b_1 - b_2) + (b_2 - b_3) + (b_3 - b_4) = b_1 - b_4 \\
    &\vdots \\
    S_N &= (b_1 - b_2) + (b_2 - b_3) + \dots + (b_N - b_{N+1}) = b_1 - b_{N+1}
\end{align*}
Luego, haciendo tender $N\to \infty$ se tiene:
\[ \displaystyle\lim_{N\to\infty}S_N = \lim_{N\to\infty}(b_1 - b_{N+1}) = b_1 - \lim_{N\to\infty}b_{N+1}. \]

\begin{definition}{Serie telescópica}{def_serie_telescopica}
    Una serie de la forma $\displaystyle\sum_{n=1}^{\infty}(b_n - b_{n+1})$ se llama \textbf{serie telescópica}, y cumple que:
    \begin{itemize}
        \item Si $\{b_n\}_{n \in \mathbb{N}}$ converge, entonces la serie converge y $\displaystyle\sum_{n=1}^{\infty}(b_n - b_{n+1}) = b_1 - \lim_{n\to\infty}b_n$.
        \item Si $\{b_n\}_{n \in \mathbb{N}}$ diverge, entonces $\displaystyle\sum_{n=1}^{\infty}(b_n - b_{n+1})$ diverge.
    \end{itemize}
\end{definition}

Habitualmente para este tipo de series hay que realizar descomposiciones por fracciones parciales para obtener la forma telescópica.

\begin{example}{Convergencia de serie telescópica básica}{ej:telescopica_1}
    Analice la convergencia de $\displaystyle\sum_{n=1}^{\infty}\frac{1}{n(n+1)}$.
\end{example}

\begin{soln}
    Notemos que $\displaystyle\sum_{n=1}^{\infty}\frac{1}{n(n+1)} = \sum_{n=1}^{\infty}\left(\frac{1}{n} - \frac{1}{n+1}\right)$. Luego algunas sumas parciales corresponden a:
    \begin{align*}
        S_1 &= 1 - \frac{1}{2} \\
        S_2 &= 1 - \frac{1}{2} + \frac{1}{2} - \frac{1}{3} = 1 - \frac{1}{3} \\
        S_3 &= 1 - \frac{1}{2} + \frac{1}{2} - \frac{1}{3} + \frac{1}{3} - \frac{1}{4} = 1 - \frac{1}{4} \\
        &\vdots \\
        S_n &= 1 - \frac{1}{n+1}
    \end{align*}
    Como $\displaystyle\lim_{n\to\infty}S_n = 1 - \lim_{n\to\infty}\frac{1}{n+1} = 1 \Longrightarrow \sum_{n=1}^{\infty}\frac{1}{n(n+1)} = 1$.
\end{soln}

\begin{example}{Telescópica por fracciones parciales}{ej:telescopica_2}
    Considere la serie $\displaystyle\sum_{n=1}^{\infty}\frac{1}{n^2+3n+2}$.
    \begin{enumerate}
        \item Demuestre que la serie es de tipo telescópica.
        \item Defina la sucesión de sumas parciales.
        \item Analice su convergencia e indique su suma si corresponde.
    \end{enumerate}
\end{example}

\begin{soln}
    $\left. \right.$
    \begin{enumerate}
        \item Utilizando fracciones parciales:
        \[ \frac{1}{n^2+3n+2} = \frac{1}{(n+1)(n+2)} = \frac{A}{n+1} + \frac{B}{n+2} = \frac{A(n+2) + B(n+1)}{(n+1)(n+2)} \implies 1 = A(n+2) + B(n+1) \]
        Si $n = -1 \implies A = 1$. Si $n = -2 \implies B = -1$. Por lo tanto, la serie es telescópica y su forma correspondiente es:
        \[ \displaystyle\sum_{n=1}^{\infty}\frac{1}{n^2+3n+2} = \sum_{n=1}^{\infty}\left(\frac{1}{n+1} - \frac{1}{n+2}\right). \]
        
        \item Sumas parciales:
        \begin{align*}
            S_1 &= \frac{1}{2} - \frac{1}{3} \\
            S_2 &= \left(\frac{1}{2} - \frac{1}{3}\right) + \left(\frac{1}{3} - \frac{1}{4}\right) = \frac{1}{2} - \frac{1}{4} \\
            S_3 &= \left(\frac{1}{2} - \frac{1}{3}\right) + \left(\frac{1}{3} - \frac{1}{4}\right) + \left(\frac{1}{4} - \frac{1}{5}\right) = \frac{1}{2} - \frac{1}{5} \\
            &\vdots \\
            S_n &= \frac{1}{2} - \frac{1}{n+2}
        \end{align*}
        
        \item Calculamos el límite de la suma parcial $S_n$:
        \[ \lim_{n\to\infty}S_n = \frac{1}{2} - \lim_{n\to\infty}\frac{1}{n+2} = \frac{1}{2}. \]
        Por lo tanto, la serie converge y además su suma corresponde a:
        \[ \displaystyle \sum_{n=1}^{\infty}\frac{1}{n^2+3n+2} = \frac{1}{2}. \]
    \end{enumerate}
\end{soln}

\subsubsection*{Serie geométrica}
Es de la forma: $\displaystyle\sum_{n=0}^{\infty}ar^n$ con $r \in \mathbb{R}$. Al número $r$ lo llamamos \textbf{razón} de la serie. Las sumas parciales corresponden a:
\[ S_n = a + ar + ar^2 + ar^3 + \dots + ar^{n-1} + ar^n \]
Multiplicando por $r$:
\[ rS_n = ra + r^2a + r^3a + \dots + r^n a + r^{n+1}a \]
Restando ambas expresiones y factorizando por $S_n$:
\[ S_n(1-r) = a - ar^{n+1} \implies S_n = a\left(\frac{1-r^{n+1}}{1-r}\right) \]
Luego, haciendo tender $n\to \infty$ se tiene:
\[ \displaystyle\lim_{n\to\infty}S_n = \lim_{n\to\infty}a\left(\frac{1-r^{n+1}}{1-r}\right). \]

\begin{definition}{Serie geométrica}{def_serie_geometrica}
    Una serie de la forma $\displaystyle\sum_{n=0}^{\infty}ar^n$ se llama \textbf{serie geométrica}, y cumple que:
    \begin{itemize}
        \item Si $|r| < 1$, entonces la serie converge y $\displaystyle\sum_{n=0}^{\infty}ar^n = \frac{a}{1-r}$.
        \item Si $|r| \ge 1$, entonces la serie diverge.
    \end{itemize}
\end{definition}

\begin{example}{Cálculo de series geométricas}{ej:geometrica_1}
    Analice la convergencia de las siguientes series. Calcule su valor de ser posible:
    \begin{multicols}{2}
    \begin{enumerate}
        \item $\displaystyle\sum_{n=0}^{\infty}2\left(\frac{2}{5}\right)^n$,
        \item $\displaystyle\sum_{n=1}^{\infty}\frac{3^{1-2n}}{2^n}$.
    \end{enumerate}
    \end{multicols}
\end{example}

\begin{soln}
    $\left. \right.$
    \begin{enumerate}
        \item Se tiene que $|r| = \left|\frac{2}{5}\right| = \frac{2}{5} < 1$, luego la serie converge y:
        \[ \sum_{n=0}^{\infty}2\left(\frac{2}{5}\right)^n = \frac{2}{1-\frac{2}{5}} = \frac{2}{\frac{3}{5}} = \frac{10}{3} \]
        
        \item Notemos que:
        \[ \sum_{n=1}^{\infty}\frac{3^{1-2n}}{2^n} = \sum_{n=1}^{\infty}\frac{3 \cdot 3^{-2n}}{2^n} = \sum_{n=1}^{\infty}3\left(\frac{1}{9^n \cdot 2^n}\right) = \sum_{n=1}^{\infty}3\left(\frac{1}{18^n}\right) = \sum_{n=1}^{\infty}3\left(\frac{1}{18}\right)^n \]
        Pero fijarse en el índice, comienza desde $1$ y debe ser desde $0$. Debemos acomodar la serie para que el índice comience desde $0$:
        \[ \sum_{n=1}^{\infty}3\left(\frac{1}{18}\right)^n = \left[3 + \sum_{n=1}^{\infty}3\left(\frac{1}{18}\right)^n\right] - 3 = \sum_{n=0}^{\infty}3\left(\frac{1}{18}\right)^{n+1} - 3 = \frac{3}{1-\frac{1}{18}} - 3 = \frac{3}{17} \]
        Dado que $|r| = \frac{1}{18} < 1$, la serie converge y su suma es $\frac{3}{17}$.
    \end{enumerate}
\end{soln}

\begin{remark}{Reindexar una serie}{obs_reindexar}
    Reindexar una serie es una técnica que consiste en cambiar el índice de la serie para facilitar su análisis o cálculo. Es posible siempre y cuando se preserve el orden de sus términos. No cambia la convergencia, ni su valor:
    \[ \sum_{n=k}^{\infty}a_n = \sum_{n=k+h}^{\infty}a_{n-h} = \sum_{n=k-h}^{\infty}a_{n+h} \]
    En el ejemplo anterior:
    \[ \sum_{n=1}^{\infty}3\left(\frac{1}{18}\right)^n = \sum_{n=0}^{\infty}3\left(\frac{1}{18}\right)^{n+1} = \sum_{n=0}^{\infty}3\left(\frac{1}{18}\right)\left(\frac{1}{18}\right)^n = \frac{\frac{1}{6}}{1-\frac{1}{18}} = \frac{3}{17} \]
\end{remark}

\subsection{Criterios de convergencia para series de términos no negativos}

Los siguientes criterios de convergencia se aplican a series de términos no negativos, es decir, a series que contienen términos positivos o nulos.

\subsubsection{Criterio de la integral}

\begin{theorem}{Criterio de la integral}{criterio_integral}
    Sea $\{ a_{n} \}_{n \in \mathbb{N}}$ una sucesión de términos positivos y suponga que $a_{n}=f(n)$, donde $f$ es una función decreciente, positiva y continua en $[1, \infty)$. Entonces:
    \begin{equation}
        \sum_{n=1}^{\infty}a_{n} \text{ converge} \Longleftrightarrow \int_{1}^{\infty}f(x)dx \text{ converge}
    \end{equation}
\end{theorem}

\begin{example}{Aplicación del criterio de la integral}{ej:criterio_integral_1}
    Analice con el criterio de la integral la convergencia de las siguientes series:
    \begin{multicols}{2}
    \begin{enumerate}
        \item $\displaystyle \sum_{n=1}^{\infty}e^{-n}$,
        \item $\displaystyle \sum_{n=1}^{\infty}\frac{e^{\arctan(n)}}{n^{2}+1}$.
    \end{enumerate}
    \end{multicols}
\end{example}

\begin{soln}
    $\left. \right.$
    \begin{enumerate}
        \item Para la serie $\displaystyle \sum_{n=1}^{\infty}e^{-n}$ se define $f(x):=e^{-x}$ la cual es continua, decreciente y positiva en $[1, \infty)$. Además:
        \[ \int_{1}^{\infty}e^{-x}dx = \lim_{b\to\infty}\int_{1}^{b}e^{-x}dx = \lim_{b\to\infty} -e^{-x}\Big|_{1}^{b} = \lim_{b\to\infty} -(e^{-b}-e^{-1}) = e^{-1} \]
        Como $\displaystyle \int_{1}^{\infty}e^{-x}dx$ converge, entonces $\displaystyle \sum_{n=1}^{\infty}e^{-n}$ también converge.
        
        \item Sea $f(x) = \dfrac{e^{\arctan(x)}}{x^{2}+1}$, la cual es una función continua, positiva y decreciente pues
        \[ f^{\prime}(x) = \frac{e^{\arctan(x)}\cdot\frac{1}{x^{2}+1}\cdot(x^{2}+1) - e^{\arctan(x)}\cdot 2x}{(x^{2}+1)^{2}} = -\frac{e^{\arctan(x)}}{(x^{2}+1)^{2}}(2x-1) < 0 \]
        Entonces:
        \[ \int_{1}^{\infty}\frac{e^{\arctan(x)}}{x^{2}+1}dx = \lim_{b\to\infty}\int_{1}^{b}\frac{e^{\arctan(x)}}{x^{2}+1}dx \]
        Haciendo $u=\arctan(x)$ y $du=\dfrac{1}{x^{2}+1}dx$, con los límites $x=b \Rightarrow u=\arctan(b)$ y $x=1 \Rightarrow u=\dfrac{\pi}{4}$:
        \[ \lim_{b\to\infty}\int_{\frac{\pi}{4}}^{\arctan(b)}e^{u}du = \lim_{b\to\infty}e^{u}\Big|_{\frac{\pi}{4}}^{\arctan(b)} = \lim_{b\to\infty}(e^{\arctan(b)} - e^{\frac{\pi}{4}}) = e^{\frac{\pi}{2}} - e^{\frac{\pi}{4}} \]
        Como la integral converge, entonces por el criterio de la integral $\displaystyle \sum_{n=1}^{\infty}\frac{e^{\arctan(n)}}{n^{2}+1}$ también converge.
    \end{enumerate}
\end{soln}

\begin{remark}{Sobre el valor de convergencia y la Serie $p$}{obs_criterio_integral}
    \begin{itemize}
        \item La serie y la integral no necesariamente convergen al mismo valor.
        \item \textbf{(Serie -p)}: Utilizando el criterio de la integral, se puede probar que
        \[ \sum_{n=1}^{\infty}\frac{1}{n^{p}} = \begin{cases} \text{ Converge} & ,\text{si } p > 1 \\ & \\ \text{ Diverge} & ,\text{si } p \le 1 \end{cases}.\]
        Este tipo de series será muy útil para comparar con otras series y así determinar su convergencia o divergencia.
    \end{itemize}
\end{remark}

\begin{problem}{}{ej:criterio_integral_2}
    Analice la convergencia de las siguientes series utilizando el criterio de la integral:
    \begin{multicols}{2}
        \begin{enumerate}
            \item $\displaystyle \sum_{n=1}^{\infty}\frac{1}{\sqrt{n}}$
            \item $\displaystyle \sum_{n=1}^{\infty}\frac{1}{n^{2}+1}$
        \end{enumerate}
    \end{multicols}
\end{problem}

\subsubsection{Criterio de comparación directa}

\begin{theorem}{Criterio de comparación directa}{criterio_comparacion_directa}
    Sean $\{ a_{n} \}_{n \in \N}$, $\{ b_{n} \}_{n \in \N}\in \mathbb{R}$ dos sucesiones de términos no negativos tales que $\forall n \ge N : 0 \le a_{n} \le b_{n}$, para algún $N$. Se cumple que:
    \begin{itemize}
        \item Si $\displaystyle \sum_{n=1}^{\infty}b_{n}$ converge, entonces $\displaystyle \sum_{n=1}^{\infty}a_{n}$ converge.
        \item Si $\displaystyle \sum_{n=1}^{\infty}a_{n}$ diverge, entonces $\displaystyle \sum_{n=1}^{\infty}b_{n}$ diverge.
    \end{itemize}
\end{theorem}

\begin{example}{Uso de cotas superiores e inferiores}{ej:criterio_comparacion_directa_1}
    Estudie la convergencia de las siguientes series:
    \begin{multicols}{2}
    \begin{enumerate}
        \item $\displaystyle \sum_{n=1}^{\infty}\frac{\ln(n)}{n}$,
        \item $\displaystyle \sum_{n=1}^{\infty}\frac{1}{1+6^{n+1}}$.
    \end{enumerate}
    \end{multicols}
\end{example}

\begin{soln}
    $\left. \right.$
    \begin{enumerate}
        \item Primero se observa que: $1 < \ln(n)$, $\forall n \ge 3 \implies \displaystyle \sum_{n=3}^{\infty}\frac{1}{n} < \displaystyle \sum_{n=3}^{\infty}\frac{\ln(n)}{n}$. Además, notemos que hay que restar los primeros términos de la serie $\displaystyle \sum_{n=1}^{\infty}\frac{1}{n}$ para que la comparación sea válida, es decir:
        \[ \sum_{n=3}^{\infty}\frac{1}{n} = \sum_{n=1}^{\infty}\frac{1}{n} - \sum_{n=1}^{2}\frac{1}{n}, \]
        y como $\displaystyle \sum_{n=1}^{\infty}\frac{1}{n}$ es una Serie-p con $p=1$, entonces diverge y por criterio de comparación directa:
        \[ \sum_{n=1}^{\infty}\frac{\ln(n)}{n} = \sum_{n=1}^{2}\frac{\ln(n)}{n} + \sum_{n=3}^{\infty}\frac{\ln(n)}{n} = \frac{\ln(2)}{2} + \sum_{n=3}^{\infty}\frac{\ln(n)}{n} \]
        diverge.
        
        \item Vemos que:
        \[ \frac{1}{1+6^{n+1}} < \frac{1}{6^{n+1}} = \frac{1}{6}\left(\frac{1}{6}\right)^{n}, \quad \forall n \ge 1 \implies \sum_{n=1}^{\infty}\frac{1}{1+6^{n+1}} < \frac{1}{6}\sum_{n=1}^{\infty}\left(\frac{1}{6}\right)^{n}. \]
        Además $\displaystyle \sum_{n=1}^{\infty}\left(\frac{1}{6}\right)^{n} = \displaystyle \sum_{n=0}^{\infty}\left(\frac{1}{6}\right)^{n} - 1$ y la serie $\displaystyle \sum_{n=0}^{\infty}\left(\frac{1}{6}\right)^{n}$ converge por ser serie geométrica con $\abs{r}=\dfrac{1}{6}<1$. Luego $\displaystyle \sum_{n=1}^{\infty}\left(\frac{1}{6}\right)^{n}$ converge y por comparación directa $\displaystyle \sum_{n=1}^{\infty}\frac{1}{1+6^{n+1}}$ también converge.
    \end{enumerate}
\end{soln}

\begin{problem}{Ejercicios de comparación directa}{ej:criterio_comparacion_directa_3}
    Decida si las siguientes series convergen utilizando el criterio de comparación directa:
    \begin{multicols}{2}
        \begin{enumerate}
            \item $\displaystyle\sum_{n=1}^{\infty}\frac{1}{n^{4}+\sqrt{n}+1}$
            \item $\displaystyle\sum_{n=1}^{\infty}\frac{6^{n}}{n8^{n}+\sqrt{n}}$
            \item $\displaystyle\sum_{n=1}^{\infty}\frac{7}{\sqrt{n}-\frac{1}{2}}$
            \item $\displaystyle\sum_{n=1}^{\infty}\frac{2\sin(n!)}{2n^{2}+1}$
        \end{enumerate}
    \end{multicols}
\end{problem}

\subsubsection{Criterio de comparación por paso al límite}

\begin{theorem}{Criterio de comparación por paso al límite}{criterio_comparacion_limite}
    Sean $\{ a_{n} \}_{n \in \N}$, $\{ b_{n} \}_{n \in \N}$ dos sucesiones de términos positivos tales que $\displaystyle \lim_{n \to \infty} \frac{a_n}{b_n} = L$. Se cumple que:
    \begin{itemize}
        \item Si $L \in \R^{+}$, entonces $\displaystyle \sum_{n=1}^{\infty}a_{n}$ y $\displaystyle \sum_{n=1}^{\infty}b_{n}$ convergen, o bien, ambas series divergen.
        \item Si $L = 0$ y $\displaystyle \sum_{n=1}^{\infty}b_{n}$ converge, entonces $\displaystyle \sum_{n=1}^{\infty}a_{n}$ converge.
        \item Si $L = +\infty$ y $\displaystyle \sum_{n=1}^{\infty}b_{n}$ diverge, entonces $\displaystyle \sum_{n=1}^{\infty}a_{n}$ diverge.
    \end{itemize}
\end{theorem}

\begin{example}{Comportamiento asintótico de términos generales}{ej:criterio_comparacion_limite_1}
    Analice la convergencia de las siguientes series mediante comparación por paso al límite:
    \begin{multicols}{2}
    \begin{enumerate}
        \item $\displaystyle \sum_{n=1}^{\infty}{\frac{\sqrt{n}\, 4^{2n}}{5n-3}}$,
        \item $\displaystyle \sum_{n=1}^{\infty}\frac{\sqrt[3]{n}}{4n^2 +3n+2}$.
    \end{enumerate}
    \end{multicols}
\end{example}

\begin{soln}
    $\left. \right.$
    \begin{enumerate}
        \item Consideremos los términos generales $a_n = \dfrac{\sqrt{n}\, 4^{2n}}{5n-3}$ y $b_n = \frac{1}{\sqrt{n}}$. Se tiene el límite:
        \[ L = \lim_{n \to \infty} \frac{a_n}{b_n} = \lim_{n \to \infty} \frac{\dfrac{\sqrt{n}\, 4^{2n}}{5n-3}}{\sqrt{n}} = \lim_{n \to \infty} \frac{n\, 4^{2n}}{5n-3} = \lim_{n \to \infty} \frac{4^{2n}}{\l 5-\frac{3}{n} \r} = +\infty. \]
        Como $\displaystyle \sum_{n=1}^{\infty} \frac{1}{\sqrt{n}}$ es una serie-p con $p = \frac{1}{2} < 1$, entonces diverge. Por tanto, por el criterio de comparación por paso al límite, la serie original también diverge.
        
        \item Utilizando los términos generales $a_n = \dfrac{\sqrt[3]{n}}{4n^2 +3n+2}$ y $b_n = \dfrac{1}{n^{5/3}}$. Se tiene el límite:
        \[ L = \lim_{n \to \infty} \frac{a_n}{b_n} = \lim_{n \to \infty} \frac{\dfrac{\sqrt[3]{n}}{4n^2 +3n+2}}{\dfrac{1}{n^{5/3}}} = \lim_{n \to \infty} \frac{n^{5/3} \cdot \sqrt[3]{n}}{4n^2 +3n+2} = \lim_{n \to \infty} \frac{n^2}{4n^2 +3n+2} = \frac{1}{4}. \]
        Como $\displaystyle \sum_{n=1}^{\infty} \frac{1}{n^{5/3}}$ es una serie-p con $p = \frac{5}{3} > 1$, entonces converge. Por tanto, por el criterio de comparación por paso al límite, la serie original también converge.
    \end{enumerate}
\end{soln}

\begin{problem}{Ejercicios de comparación al límite}{ej:criterio_comparacion_limite_3}
    Analice si las siguientes series convergen o divergen:
    \begin{multicols}{2}
        \begin{enumerate}
            \item $\displaystyle\sum_{n=1}^{\infty}\frac{4n-3}{n^3+3n^2+1}$
            \item $\displaystyle\sum_{n=1}^{\infty}\frac{5n}{\sqrt{2n^4+1}}$
            \item $\displaystyle\sum_{n=1}^{\infty}\sin{\l \dfrac{\pi}{n} \r}$
            \item $\displaystyle\sum_{n=1}^{\infty}\frac{n^2}{n^5 +\sqrt{n}}$
        \end{enumerate}
    \end{multicols}
\end{problem}

\subsection{Criterios de convergencia para series de términos de signo variable}

Los siguientes criterios se aplican a series que pueden contener términos negativos o que cambian de signo dinámicamente, tales como las \textbf{series alternantes} de la forma $\displaystyle \sum_{n=1}^{\infty}(-1)^n a_n$ o $\displaystyle \sum_{n=1}^{\infty}(-1)^{n+1} a_n$.

\subsubsection{Criterio de Leibniz}

\begin{theorem}{Criterio de Leibniz}{criterio_leibniz}
    Sea $\{ a_{n} \}_{n \in \N}$ una sucesión de términos positivos tales que:
    \begin{itemize}
        \item $a_n \geq a_{n+1}$ para todo $n \in \N$,
        \item $\displaystyle \lim_{n \to \infty} a_n = 0$.
    \end{itemize}
    Entonces, las series $\displaystyle \sum_{n=1}^{\infty}(-1)^n a_n$ y $\displaystyle \sum_{n=1}^{\infty}(-1)^{n+1} a_n$ son convergentes.
\end{theorem}

\begin{example}{Análisis de series alternantes}{ej:criterio_leibniz_1}
    Estudie la convergencia de las siguientes series alternantes:
    \begin{multicols}{2}
    \begin{enumerate}
        \item $\displaystyle \sum_{n=1}^{\infty}\frac{(-1)^{n}}{\sqrt{n}}$,
        \item $\displaystyle \sum_{n=1}^{\infty}{(-1)^{n+1} \frac{3n+2}{4n^2-3}}$.
    \end{enumerate}
    \end{multicols}
\end{example}

\begin{soln}
    $\left. \right.$
    \begin{enumerate}
        \item Debemos verificar las condiciones sobre $a_n = \dfrac{1}{\sqrt{n}}$:
        \begin{itemize}
            \item Definimos $f(n) := a_n = \dfrac{1}{\sqrt{n}}$, la cual es una función positiva y decreciente, pues $f^{\prime}(n) = -\dfrac{1}{2n^{3/2}} < 0$ para todo $n \in \N$.
            \item $\displaystyle \lim_{n \to \infty} a_n = \displaystyle \lim_{n \to \infty} \frac{1}{\sqrt{n}} = 0$.
        \end{itemize}
        Como se cumplen ambas, por Leibniz, la serie alternante converge.
        
        \item Sea el término general $a_n = \dfrac{3n+2}{4n^2-3}$. Verificamos:
        \begin{itemize}
            \item Definimos $f(n) := a_n$, que es positiva y decreciente para todo $n\in\mathbb{N}$ ya que:
            \[ f^{\prime}(n) = \dfrac{3 (4n^2 -3) - 8n(3n+2)}{(4n^2-3)^2} = \dfrac{- \l 12n^2+16n+9 \r}{(4n^2-3)^2} < 0. \]
            \item $\displaystyle \lim_{n \to \infty} a_n = \displaystyle \lim_{n \to \infty} \frac{3n+2}{4n^2-3} = 0$.
        \end{itemize}
        Cumpliéndose las condiciones, la serie converge por el criterio de Leibniz.
    \end{enumerate}
\end{soln}

\begin{problem}{Ejercicios de series alternantes}{}
    Estudie la convergencia de las siguientes series alternantes:
    \begin{multicols}{2}
        \begin{enumerate}
            \item $\displaystyle \sum_{n=1}^{\infty}(-1)^{n}\frac{n+2}{5^n}$
            \item $\displaystyle \sum_{n=1}^{\infty}\frac{\cos{\l \pi n \r} \l n+2 \r}{\sqrt{n^2 +n}}$
            \item $\displaystyle \sum_{n=1}^{\infty}(-1)^{n} \frac{n^2}{n^2+1}$
            \item $\displaystyle \sum_{n=1}^{\infty}(-1)^{n} \frac{2^n \, \l \frac{1}{2} \r^n}{2n-1}$
        \end{enumerate}
    \end{multicols}
\end{problem}

\subsubsection{Criterio de la razón y de la raíz}

\begin{theorem}{Criterio de la razón}{criterio_razon}
    Sea $\displaystyle \sum_{n=1}^{\infty} a_n$ una serie y suponga que $r = \displaystyle \lim_{n \to \infty} \abs{\frac{a_{n+1}}{a_n}}$ existe. Entonces:
    \begin{itemize}
        \item Si $r < 1$, entonces la serie es convergente.
        \item Si $r > 1$, entonces la serie es divergente.
        \item Para $r = 1$ el criterio no proporciona información.
    \end{itemize}
\end{theorem}

\begin{example}{Series con factoriales y potencias}{ej:criterio_razon_1}
    Estudie la convergencia de las siguientes series utilizando el criterio de la razón:
    \begin{multicols}{2}
    \begin{enumerate}
        \item $\displaystyle \sum_{n=1}^{\infty} \frac{n!}{n^n}$,
        \item $\displaystyle \sum_{n=1}^{\infty}(-1)^n \frac{n!}{2^n}$.
    \end{enumerate}
    \end{multicols}
\end{example}

\begin{soln}
    $\left. \right.$
    \begin{enumerate}
        \item Calculamos el límite de la razón:
        \[ r = \lim_{n \to \infty} \left| \frac{a_{n+1}}{a_n} \right| = \lim_{n \to \infty} \left| \frac{(n+1)!}{(n+1)^{n+1}} \cdot \frac{n^n}{n!} \right| = \lim_{n \to \infty} \left( \frac{n}{n+1} \right)^n = \frac{1}{e} < 1. \]
        Como $r = \frac{1}{e} < 1$, por el criterio de la razón, la serie converge.
        
        \item Calculamos el límite del cociente en valor absoluto:
        \[ r = \lim_{n \to \infty} \left| \frac{(-1)^{n+1} (n+1)!}{2^{n+1}} \cdot \frac{2^n}{(-1)^n n!} \right| = \lim_{n \to \infty} \frac{n+1}{2} = +\infty > 1. \]
        Como $r = +\infty > 1$, por el criterio de la razón, la serie diverge.
    \end{enumerate}
\end{soln}

\begin{remark}{Casos no concluyentes ($r=1$)}{obs_criterio_razon}
    El criterio de la razón no proporciona información para $r=1$. Por ejemplo, la serie $\displaystyle \sum_{n=1}^{\infty} \frac{1}{n}$ verifica $r = 1$ y diverge, mientras que la serie $\displaystyle \sum_{n=1}^{\infty} \frac{1}{n^2}$ también verifica $r = 1$, pero converge. Cuando ocurre esto, es obligatorio recurrir a otros criterios alternativos.
\end{remark}

\begin{problem}{Ejercicios del criterio de la razón}{ej:criterio_razon_4}
    Analice la convergencia de las siguientes series mediante el criterio de la razón:
    \begin{multicols}{2}
        \begin{enumerate}
            \item $\displaystyle \sum_{n=1}^{\infty} \frac{1}{n!}$
            \item $\displaystyle \sum_{n=0}^{\infty} \frac{n!}{e^{n+1}}$
            \item $\displaystyle \sum_{n=1}^{\infty} \frac{n^2}{3^{n+2}}$
            \item $\displaystyle \sum_{n=1}^{\infty} \frac{(-1)^{n+1}}{n \, 5^n}$
        \end{enumerate}
    \end{multicols}
\end{problem}

\begin{theorem}{Criterio de la raíz}{criterio_raiz}
    Sea $\displaystyle \sum_{n=1}^{\infty} a_n$ una serie y suponga que $r = \displaystyle \lim_{n \to \infty} \sqrt[n]{|a_n|}$ existe. Entonces:
    \begin{itemize}
        \item Si $r < 1$, entonces la serie es convergente.
        \item Si $r > 1$, entonces la serie es divergente.
        \item Para $r = 1$ el criterio no proporciona información.
    \end{itemize}
\end{theorem}

\begin{example}{}{ej:criterio_raiz_1}
    Analice si las siguientes series convergen o divergen utilizando el criterio de la raíz:
    \begin{multicols}{2}
    \begin{enumerate}
        \item $\displaystyle \sum_{n=1}^{\infty} \l \frac{\ln{(n)}}{n} \r^n$,
        \item $\displaystyle \sum_{n=1}^{\infty} \frac{2^n}{n^n}$.
    \end{enumerate}
    \end{multicols}
\end{example}

\begin{soln}
    $\left. \right.$
    \begin{enumerate}
        \item Calculamos el límite aplicando la raíz n-ésima:
        \[ r = \lim_{n \to \infty} \sqrt[n]{\left|\frac{\ln(n)}{n}\right|^n} = \lim_{n \to \infty} \frac{\ln(n)}{n} \overset{L'h}{=} \lim_{n \to \infty} \frac{1}{n} = 0 < 1. \]
        Como $r = 0 < 1$, por el criterio de la raíz, la serie converge.
        
        \item Calculamos el límite del radical n-ésimo:
        \[ r = \lim_{n \to \infty} \sqrt[n]{\left|\frac{2^n}{n^n}\right|} = \lim_{n \to \infty} \frac{2}{n} = 0 < 1. \]
        Como $r = 0 < 1$, por el criterio de la raíz, la serie converge.
    \end{enumerate}
\end{soln}

\begin{problem}{Ejercicios del criterio de la raíz}{ej:criterio_raiz_4}
    Utilizando el criterio de la raíz, analice la convergencia de las siguientes series:
    \begin{multicols}{2}
        \begin{enumerate}
            \item $\displaystyle \sum_{n=1}^{\infty} \l \frac{2n}{n+1} \r^n$
            \item $\displaystyle \sum_{n=1}^{\infty} \frac{n \, 3^n}{n!}$
            \item $\displaystyle \sum_{n=1}^{\infty} \frac{(-1)^n \, n^n}{e^n}$
            \item $\displaystyle \sum_{n=1}^{\infty} \frac{n^3}{3^n}$
        \end{enumerate}
    \end{multicols}
\end{problem}

\subsection{Convergencia absoluta y condicional}

\begin{definition}{Convergencia absoluta y condicional}{convergencia_absoluta_condicional}
    Sea $\displaystyle \sum_{n=1}^{\infty}a_n$ una serie de números reales. Se dice que:
    \begin{itemize}
        \item Si la serie de los valores absolutos $\displaystyle \sum_{n=1}^{\infty}|a_n| $ converge, entonces la serie original es \textbf{absolutamente convergente}.
        \item Si la serie $\displaystyle \sum_{n=1}^{\infty}|a_n|$ diverge, pero la serie $\displaystyle \sum_{n=1}^{\infty}a_n$ converge, entonces es \textbf{condicionalmente convergente}.
    \end{itemize}
\end{definition}

\begin{example}{}{ej:convergencia_absoluta_condicional_1}
    Verifique qué tipo de convergencia posee cada una de las siguientes series:
    \begin{multicols}{2}
    \begin{enumerate}
        \item $\displaystyle \sum_{n=1}^{\infty}{\frac{(-1)^n}{n^2}}$,
        \item $\displaystyle \sum_{n=1}^{\infty} \frac{(-1)^n}{n}$.
    \end{enumerate}
    \end{multicols}
\end{example}

\begin{soln}
    $\left. \right.$
    \begin{enumerate}
        \item Aplicando módulo: $\displaystyle \sum_{n=1}^{\infty}\left|\frac{(-1)^n}{n^2}\right| = \displaystyle \sum_{n=1}^{\infty} \frac{1}{n^2}$, que es una serie-p con $p=2 > 1$ (converge). Por ende, la serie es absolutamente convergente.
        
        \item Aplicando módulo: $\displaystyle \sum_{n=1}^{\infty}\left|\frac{(-1)^n}{n}\right| = \displaystyle \sum_{n=1}^{\infty} \frac{1}{n}$, que es la serie armónica ($p=1$), la cual diverge. Sin embargo, por el criterio de Leibniz se sabe que la serie alternante original sí converge. Por lo tanto, es condicionalmente convergente.
    \end{enumerate}
\end{soln}

\begin{problem}{}{}
    Analice si cada serie siguiente es absolutamente convergente, condicionalmente convergente o divergente:
    \begin{multicols}{2}
        \begin{enumerate}
            \item $\displaystyle \sum_{n=1}^{\infty} \frac{(-1)^n n^2}{n^2+1}$
            \item $\displaystyle \sum_{n=1}^{\infty} \frac{2(-1)^n}{e^n - e^{-n}}$
        \end{enumerate}
    \end{multicols}
\end{problem}

\newpage
\subsection{Tabla resumen de criterios de convergencia para series}

\begin{table}[H]
\centering
\resizebox{\textwidth}{!}{%
\renewcommand{\arraystretch}{1.5}
\begin{tabular}{|l|c|l|l|}
\hline 
\textbf{Nombre} & \textbf{Serie} & \textbf{Criterio} & \textbf{Observación} \\ \hline 
Término n-ésimo & $\displaystyle\sum_{n=1}^{\infty}a_n$ & $\displaystyle\lim_{n \rightarrow \infty}a_n \neq 0 \, \Longrightarrow \,$ Div. & No sirve para convergencia \\ \hline
Serie geométrica & $\displaystyle\sum_{n=0}^{\infty}ar^{n}$ & $\begin{cases} |r| < 1 \Longrightarrow \text{Conv. a } S=\frac{a}{1-r} \\ |r| \geq 1 \Longrightarrow \text{Div.} \end{cases}$ & Útil para comparación \\ \hline
Serie p & $\displaystyle\sum_{n=1}^{\infty}\frac{1}{n^p}$ & $\begin{cases} p>1 \Longrightarrow \text{Conv.} \\ p \leq 1 \Longrightarrow \text{Div.} \end{cases}$ & Útil para comparación \\ \hline
Integral & $\begin{cases} \displaystyle\sum_{n=1}^{\infty}a_n \\ a_n =f(n) \end{cases}$ & $\begin{cases} \displaystyle\int_{1}^{\infty}f(x) \, dx \text{ Conv. } \Longrightarrow \text{ Conv.} \\ \displaystyle\int_{1}^{\infty}f(x) \, dx \text{ Div. } \Longrightarrow \text{ Div.} \end{cases}$ & $\begin{cases} f \text{ continua, positiva,} \\ \text{decreciente y de fácil int.} \end{cases}$ \\ \hline
Comp. directa & $\begin{cases} \displaystyle\sum_{n=1}^{\infty}a_n, \, \displaystyle\sum_{n=1}^{\infty}b_n \\ 0 \leq a_n \leq b_n \end{cases}$ & $\begin{cases} \displaystyle\sum_{n=1}^{\infty}b_n \text{ Conv. } \Longrightarrow \displaystyle\sum_{n=1}^{\infty}a_n \text{ Conv.} \\ \displaystyle\sum_{n=1}^{\infty}a_n \text{ Div. } \Longrightarrow \displaystyle\sum_{n=1}^{\infty}b_n \text{ Div.} \end{cases}$ & $\begin{cases} \text{Usar series geom. o p} \\ \text{para comparar} \end{cases}$ \\ \hline
Comp. en el límite & $\displaystyle\sum_{n=1}^{\infty}a_n, \,\displaystyle\sum_{n=1}^{\infty}b_n >0$ & $\displaystyle\lim_{n \rightarrow \infty}\frac{a_n}{b_n}=L \begin{cases} L \neq 0, \text{ ambas Conv. o Div.} \\ L = 0 \text{ y } \displaystyle\sum b_n \text{ Conv. } \Longrightarrow \displaystyle\sum a_n \text{ Conv.} \\ L = \infty \text{ y } \displaystyle\sum b_n \text{ Div. } \Longrightarrow \displaystyle\sum a_n \text{ Div.}\end{cases}$ & $\begin{cases} \text{Usar series geom. o p} \\ \text{para comparar} \end{cases}$ \\ \hline
Serie Alt. (Leibniz) & $\displaystyle\sum_{n=1}^{\infty}(-1)^n a_n$ & $\begin{cases} a_n \geq a_{n+1}, \, \forall n \\ \displaystyle\lim_{n \rightarrow \infty}a_n = 0 \end{cases} \Longrightarrow \text{ Conv.}$ & $\begin{cases} \text{Si } \displaystyle\lim_{n \rightarrow \infty}a_n \neq 0 \\ \text{aplicar término n-ésimo}\end{cases}$ \\ \hline
Razón & $\displaystyle\sum_{n=1}^{\infty}a_n$ & $\displaystyle\lim_{n \rightarrow \infty}\left|\frac{a_{n+1}}{a_n}\right|= L \begin{cases} L < 1 , \text{ Conv.} \\ L > 1, \text{ Div.} \end{cases}$ & $\begin{cases} \text{No decide si } L = 1. \\ \text{Útil si } a_n \text{ tiene} \\ \text{factoriales o potencias} \end{cases}$ \\ \hline
Raíz & $\displaystyle\sum_{n=1}^{\infty}a_n$ & $\displaystyle\lim_{n \rightarrow \infty}\sqrt[n]{|a_n|}= L \begin{cases} L < 1 , \text{ Conv.} \\ L > 1, \text{ Div.} \end{cases}$ & $\begin{cases} \text{No decide si } L = 1. \\ \text{Útil si } a_n \text{ tiene potencias} \end{cases}$ \\ \hline
Serie Abs. Conv. & $\displaystyle\sum_{n=1}^{\infty}a_n$ & $\displaystyle\sum_{n=1}^{\infty}|a_n| \text{ Conv. } \Longrightarrow \displaystyle\sum_{n=1}^{\infty}a_n \text{ Abs. Conv.}$ & $\begin{cases} \text{Series de térm.} \\ \text{positivos y negativos} \end{cases}$ \\ \hline
Serie Cond. Conv. & $\displaystyle\sum_{n=1}^{\infty}a_n$ & $\displaystyle\sum_{n=1}^{\infty}a_n \text{ Conv. y } \displaystyle\sum_{n=1}^{\infty}|a_n| \text{ Div. } \Longrightarrow \displaystyle\sum_{n=1}^{\infty}a_n \text{ Cond. Conv.}$ & $\begin{cases} \text{Series de térm.} \\ \text{positivos y negativos} \end{cases}$ \\ \hline
\end{tabular}%
}
\end{table}

\newpage
\subsection{Ejercicios propuestos}

\renewcommand{\theenumi}{\arabic{enumi}}
\renewcommand{\labelenumi}{\textbf{\theenumi.-}}
\begin{enumerate}
\item Usar, si es posible, el criterio de divergencia para determinar si las siguientes series divergen.
\begin{multicols}{2}
\begin{enumerate}
\item $ \displaystyle{ \sum_{n=2}^{\infty}{ \l -1 \r^n \frac{1}{n^2} } } $
\item $ \displaystyle{ \sum_{n=1}^{\infty}{ \sin{ \l \frac{\pi}{2} - \frac{1}{n} \r} } } $
\item $ \displaystyle{ \sum_{n=1}^{\infty}{ \l 1 + \frac{1}{n} \r \ln{\l 1 + \frac{1}{n} \r }} } $
\item $ \displaystyle{ \sum_{n=0}^{\infty}{ \frac{n^2}{n + 1} } } $
\end{enumerate}
\end{multicols}

\item Estudiar la naturaleza de las siguientes series. Encontrar la suma, cuando sea posible.
\begin{multicols}{3}
\begin{enumerate}
\item $ \displaystyle{ \sum_{n=1}^{\infty}{ \l -1 \r^n \frac{3^n}{2^{n+2}} } } $
\item $ \displaystyle{ \sum_{n=1}^{\infty}{ \frac{3^{n+3}}{5^{n-1}} } } $
\item $ \displaystyle{ \sum_{n=1}^{\infty}{ \frac{2^n + 5^n}{2^n 5^n}} } $
\item $ \displaystyle{ \sum_{n=1}^{\infty}{ \frac{1}{2^{n-1}} } } $
\item $ \displaystyle{ \sum_{n=0}^{\infty}{ \frac{2^n \sqrt{3} - 3^n \sqrt{2}}{6^n} } } $
\item $ \displaystyle{26 \sum_{n=2}^{\infty}{ 3^{2-3n } } + \frac{26}{3} } $
\end{enumerate}
\end{multicols}

\item Sea $\displaystyle{a_n = \sqrt{n+1} - \sqrt{n} }$, se pide:
\begin{enumerate}
\item Calcular el $\displaystyle{ \lim_{n \rightarrow \infty}{a_n} }$.
\item Determinar si $\displaystyle{\sum_{n=1}^{\infty}{a_n}}$ es convergente.
\end{enumerate}

\item Usar, cuando sea posible, los criterios de comparación directa, comparación en el límite o de la integral para determinar si las siguientes series convergen o divergen.
\begin{multicols}{3}
\begin{enumerate}
\item $ \displaystyle{ \sum_{n=1}^{\infty}{ \frac{1}{3^n - \sin{\l n \r}} } } $
\item $ \displaystyle{ \sum_{n=2}^{\infty}{ \frac{1}{\sqrt[n]{\ln{ \l n \r }}} } } $
\item $ \displaystyle{ \sum_{n=1}^{\infty}{ \frac{\ln{ \l n+1 \r } - \ln{ \l n \r }}{\arctan{\l 2n \r}} } } $
\item $ \displaystyle{ \sum_{n=1}^{\infty}{ \frac{1}{\sqrt[3]{8n^2 - 5n}} } } $
\item $ \displaystyle{ \sum_{n=2}^{\infty}{ \frac{\l \ln{ \l n \r } \r^2}{n^3} } } $
\item $ \displaystyle{ \sum_{n=1}^{\infty}{ \frac{e^{\frac{1}{n}}}{n^2} } } $
\item $ \displaystyle{ \sum_{n=1}^{\infty}{ \frac{1}{\l n+1 \r \sqrt{\ln{\l n+1 \r}}}} } $ 
\item $ \displaystyle{ \sum_{n=1}^{\infty}{ \frac{1}{n \sqrt{n}} } } $
\item $ \displaystyle{ \sum_{n=1}^{\infty}{ \frac{n+1}{n \sqrt{n}} } } $
\item $ \displaystyle{ \sum_{n=1}^{\infty}{ \frac{1}{n 2^n} } } $
\item $ \displaystyle{ \sum_{n=1}^{\infty}{ \frac{1}{ \sqrt{n \l n+1 \r}} } } $
\item $ \displaystyle{ \sum_{n=1}^{\infty}{ \frac{2}{2^n + 3}} } $ 
\item $ \displaystyle{ \sum_{n=1}^{\infty}{ \frac{\ln{ \l n+1 \r }}{\l n+1 \r ^3} } } $
\item $ \displaystyle{ \sum_{n=1}^{\infty}{ \frac{n}{\sqrt{4n^3 +1}} } } $
\end{enumerate}
\end{multicols}

\item ¿Para qué valores de $p$ la serie $\displaystyle{\sum_{n=2}^{\infty}{ \frac{1}{n \l \ln{\l n \r} \r^p} }}$ converge?
\item ¿Para qué valores de $p$ y $q$ la serie $\displaystyle{\sum_{n=2}^{\infty}{ \frac{\l \ln{\l n \r} \r^p}{n^q} }}$ converge?

\item Usar, cuando sea posible, los criterios de la razón y de la raíz para determinar si las siguientes series convergen o divergen.
\begin{multicols}{3}
\begin{enumerate}
\item $ \displaystyle{ \sum_{n=1}^{\infty}{ \frac{n^n}{\l p+1 \r \l p+2 \r \ldots \l p+n \r } } } $
\item $ \displaystyle{ \sum_{n=1}^{\infty}{ \frac{\l n! \r^2 2^n}{\l 2n+2 \r! } } } $
\item $ \displaystyle{ \sum_{n=1}^{\infty}{2^n \frac{\sqrt[n]{n}}{n! } } } $
\item $ \displaystyle{ \sum_{n=1}^{\infty}{ \l 2n^2 + 3n +5 \r \l \frac{2}{3 } \r^n } } $
\item $ \displaystyle{ \sum_{n=1}^{\infty}{ n^2 e^{-n} } } $
\item $ \displaystyle{ \sum_{n=1}^{\infty}{ \frac{2^n n^n}{ \binom{n}{2}^n } } } $
\item $ \displaystyle{ \sum_{n=1}^{\infty}{ \frac{\l -1 \r^{n-1} n!}{6^n } } } $
\item $ \displaystyle{ \sum_{n=1}^{\infty}{ \frac{n!}{10^n} } } $
\item $ \displaystyle{ \sum_{n=1}^{\infty}{ \frac{10^n}{n!} } } $
\item $ \displaystyle{ \sum_{n=1}^{\infty}{ n \l \frac{3}{4 } \r^n } } $
\item $ \displaystyle{ \sum_{n=1}^{\infty}{ \frac{\l -1 \r^n}{n^p } } }, \,\, 0 < p <1 $
\item $ \displaystyle{ \sum_{n=1}^{\infty}{ \frac{2^n}{n^3 } } } $
\item $ \displaystyle{ \sum_{n=1}^{\infty}{ \l \frac{3}{2} \r^n \frac{n!}{n} } } $
\item $ \displaystyle{ \sum_{n=1}^{\infty}{ \frac{n}{\l n+1 \r 2^n} } } $
\end{enumerate}
\end{multicols}

\item Determinar si las siguientes series divergen, convergen condicionalmente o convergen absolutamente.
\begin{multicols}{3}
\begin{enumerate}
\item $ \displaystyle{ \sum_{n=1}^{\infty}{ \l -1 \r^n \frac{n+1}{5^n} } } $
\item $ \displaystyle{ \sum_{n=1}^{\infty}{ \l -1 \r^{n+1} \frac{2}{e^n - e^{-n}} } } $
\item $ \displaystyle{ \sum_{n=1}^{\infty}{ \l -1 \r^n \frac{n^2}{n^2 +1} } } $
\item $ \displaystyle{ \sum_{n=1}^{\infty}{ \l -1 \r^n \frac{\cos{\l n \pi \r}}{n+1} } } $
\item $ \displaystyle{ \sum_{n=1}^{\infty}{ \l -1 \r^{n+1} \frac{\ln{\l n+1 \r}}{n+1} } } $
\item $ \displaystyle{ \sum_{n=1}^{\infty}{ \l -1 \r^{n-1} \frac{\ln{\l n \r}}{n^2} } } $
\item $ \displaystyle{ \sum_{n=1}^{\infty}{ \l -1 \r^{n-1} \frac{\l \frac{4}{3} \r^n}{n^2} } } $
\item $ \displaystyle{ \sum_{n=1}^{\infty}{ \l -1 \r^n \frac{1}{n^2 \ln{\l n \r}} } } $
\end{enumerate}
\end{multicols}

\end{enumerate}
